\section[L. 海底捞月]{\myboxxii{      \BF{\TT{L}}} 海底捞月}

\textit{运行时间限制}   2 seconds

\textit{内存使用限制} 1 Gigabyte

\textit{读入文件} \qquad 标准输入

\textit{输出文件} \qquad 标准输出


\begin{quote}
	\textit{\_\_\_, \_\_\_ 面色一沉：河底捞鱼是什么意思？}
\end{quote}

\subsection{题目描述}

TSUCE 的期末考试结束了！阅卷老师大 K 看了看小 T 惨不忍睹的试卷。

因此他决定使用以下的招数，拯救考了 $x$ 分的小 T：

\begin{itemize}
	\item 【连根拔起】，这个招数至多可以使用 $a_1$ 次，将小 T 当前的分数 $x$ 开根号乘 10 下取整，也即 $x \leftarrow \lfloor 10 \sqrt{x} \rfloor$.
	\item 【三七分成】，这个招数至多可以使用 $a_2$ 次，将小 T 当前的分数 $x$ 乘上 $0.7$ 再加上 30 再下取整， 也即 $x \leftarrow \lfloor 0.7x + 30 \rfloor$.
	\item 【种五得六】，这个招数至多可以使用 $a_3$ 次，将小 T 当前的分数 $x$ 乘上 $1.2$ 再 下取整，也即 $x \leftarrow \lfloor 1.2x \rfloor$.
	\item 【五福临门】，这个招数至多可以使用 $a_4$ 次，将小 T 当前的分数 $x$ 加 $5$，也即 $x \leftarrow x + 5$.
\end{itemize}

四种招数可以按任意顺序使用，请帮大 K 求出一个合理的使用招数的方式，来让小 T 获得最高的分数。

\subsection{输入格式}

第一行输入一个整数 $T$ ($1 \le T \le 3 \times 10^5$), 表示测试数据的数量。

每组测试数据输入一行 5 个整数 $x$, $a_1$, $a_2$, $a_3$, $a_4$ ($1 \le x \le 100, 0 \le a_1, a_2, a_3, a_4 \le 400$), 表示小 T 的初始分数，以及四种招数的可用次数上限。

\subsection{输出格式}

对于每组测试数据，输出一行一个整数表示小 T 可以获得的最高分数。

\subsection{样例 1}

\begin{minipage}{\linewidth}
	\begin{minipage}[t]{0.49\linewidth}
		\begin{lstlisting}[title=\texttt{standard input}]
7
1 1 1 1 1
50 2 1 1 2
99 0 1 2 1
10 2 0 0 3
100 400 400 400 400
75 1 7 3 0
1 3 7 4 3
\end{lstlisting}
	\end{minipage}
	\begin{minipage}[t]{0.49\linewidth}
		\begin{lstlisting}[title=\texttt{standard output}]
70
118
148
71
98703086085238798121046930830507710
168
232
\end{lstlisting}
	\end{minipage}
\end{minipage}


\subsection{样例 1 解释}

对于第 $1$ 组测试数据，小 T 一开始豪取 $1$ 分，大 K 每种招数都至多能使用一次，大 K 可以采用如下的策略：

\begin{itemize}
	\item 首先用 【三七分成】，小 T 变为 $\lfloor 0.7x+30 \rfloor = \lfloor 0.7 \times 1 +30 \rfloor = 30$ 分。
	\item 随后用 【连根拔起】，小 T 变为 $\lfloor 10\sqrt{x} \rfloor = \lfloor 10\sqrt{30} \rfloor = 54$ 分。
	\item 随后用 【五福临门】，小 T 变为 $54 + 5 = 59$ 分。
	\item 随后用 【种五得六】，小 T 变为 $\lfloor 1.2x \rfloor = \lfloor 1.2 \times 59 \rfloor = 70$ 分。
\end{itemize}

% \subsection{提示}

% 请选用合适的数据类型进行运算，样例 $1$ 种第 5 组数据的输出有 \[34 < \log_{10}(98\,703\,086\,085\,238\,798\,121\,046\,930\,830\,507\,710) < 35.\]




